\section{Introduction}
\label{sec:intro}

Translating natural language into executable Structured Query Language, or Text-to-SQL, has emerged as a pivotal capability for enabling non-expert users to interrogate large-scale relational databases. Despite remarkable progress driven by Large Language Models, or LLMs, the gap between research prototypes and production-grade interfaces remains significant. Current systems face three interrelated challenges.

First, Schema Ambiguity presents a major hurdle, as enterprise databases routinely contain hundreds of tables with cryptic naming conventions, causing models to hallucinate non-existent columns or misinterpret join paths \cite{li_can_2023, floratou2024broken}. Second, a flawed Error Taxonomy limits progress because existing correction mechanisms conflate syntactical errors, which are detectable by the database engine, with logical errors, which are semantically incorrect but syntactically valid. Consequently, a unified correction prompt wastes valuable reasoning capacity on trivially diagnosable failures. Third, the Cost--Accuracy Trade-off remains problematic, as monolithic frontier-model approaches such as GPT-4 or Claude~3 achieve high accuracy but incur prohibitive per-query costs, making large-scale batch evaluation economically infeasible.

To address these challenges, we present \highlight{AgentSQL-Asym}, a state-of-the-art asymmetric Text-to-SQL framework orchestrated as a stateful \textbf{LangGraph State Machine}. Unlike sequential pipelines, AgentSQL-Asym dynamically routes queries based on complexity, decoupling schema exploration from generation, and separating semantic validation from high-reasoning syntactical correction. The framework employs an \textbf{Adaptive Router} to bypass schema enrichment for simple queries, and incorporates a \textbf{Resilient SQL Validator} loop running local SQLite sandbox executions with a \textbf{Targeted Corrector} that runs for up to three iterations to resolve failures on complex tasks. By reserving premium LLM API calls strictly for SQL generation and correction, AgentSQL-Asym maintains extreme cost-efficiency.

\medskip
\noindent\textbf{Contributions.} We make the following principal contributions in this work. We introduce a stateful, asymmetric multi-agent framework orchestrated via \textbf{LangGraph}, utilizing an Adaptive Router node to dynamically bypass offline schema exploration for simple queries to minimize latency. To ensure robust query generation, we implement a dual-candidate generation strategy in the SQL Generator node that synthesizes queries under both DDL and Markdown schema representations, selecting the optimal initial query based on sandbox execution. Furthermore, we develop a three-tier Semantic State Evaluation layer, termed the \texttt{Resilient SQL Validator}, which classifies execution tracebacks into \textsc{SyntaxError}, \textsc{EmptyResult}, and \textsc{NullResult} states, each providing targeted suggestions to the Corrector node. For high-throughput scalability, we incorporate optimizations including a round-robin \texttt{KeyRotator} that automatically rotates API keys on HTTP~429 rate limits, alongside transparent fallback models to ensure zero-sleep resilience. Finally, we conduct a comprehensive empirical evaluation on the full BIRD Mini-Dev benchmark under a highly cost-efficient asymmetric model budget, achieving an outstanding Execution Accuracy of \textbf{80.20\%}, establishing a new paradigm for practical, high-performance Text-to-SQL deployment.

